% INSTRUCTOR-ONLY ANSWER KEY: do not distribute or publish with student materials.
% Generated by generate_core_syntax_bank.py. Edit the generator, not this file.
\def\MAKEANSWERKEY{1}
\documentclass[9pt]{extarticle}

% Generated by generate_core_syntax_bank.py. Edit the generator, not this file.
\usepackage[letterpaper,margin=0.55in,includehead,includefoot,
  headheight=15pt,headsep=6pt,footskip=18pt]{geometry}
\usepackage[T1]{fontenc}
\usepackage{amsmath,amssymb}
\usepackage{microtype}
\usepackage[dvipsnames]{xcolor}
\usepackage{multicol}
\usepackage{fancyhdr}
\usepackage[hidelinks]{hyperref}

\definecolor{CoreNavy}{HTML}{18364A}
\definecolor{CoreTeal}{HTML}{176B6B}
\definecolor{CoreTint}{HTML}{EDF4F8}

\newif\ifanswers
\ifdefined\MAKEANSWERKEY
  \answerstrue
\else
  \answersfalse
\fi

\setlength{\parindent}{0pt}
\setlength{\parskip}{0pt}
\setlength{\columnsep}{0.25in}
\setlength{\columnseprule}{0.25pt}
\setlength{\emergencystretch}{1.2em}
\raggedbottom

\newcommand{\currentcore}{Core Practice}
\newcommand{\currentform}{}
\pagestyle{fancy}
\fancyhf{}
\fancyhead[L]{\textbf{COMP 2300}}
\fancyhead[C]{\textbf{\ifanswers Instructor Key\else Read --- Parse --- Judge\fi}}
\fancyhead[R]{\currentform}
\fancyfoot[L]{\ifanswers Answers, diagnoses, and minimal repairs\else Read it; parse it; then judge it.\fi}
\fancyfoot[C]{\thepage}
\fancyfoot[R]{20 items}
\renewcommand{\headrulewidth}{0.5pt}
\renewcommand{\footrulewidth}{0.3pt}

\newcounter{tfitem}
\newcommand{\CoreForm}[5]{%
  \clearpage
  \renewcommand{\currentcore}{#1 Core}%
  \renewcommand{\currentform}{Form #2}%
  \setcounter{tfitem}{0}%
  {\color{CoreNavy}\Large\bfseries #1 Core: Reading \& Interpreting --- Form #2}\par
  \vspace{2pt}\color{CoreNavy}\rule{\linewidth}{0.8pt}\color{black}\par
  \vspace{3pt}
  \ifanswers\else
    \textbf{Name:} \rule{0.42\linewidth}{0.4pt}\hfill
    \textbf{Date:} \rule{0.22\linewidth}{0.4pt}\par
    \vspace{4pt}
  \fi
  \fcolorbox{CoreTeal}{CoreTint}{\parbox{0.965\linewidth}{%
    \ifanswers\scriptsize\else\footnotesize\fi
    \textbf{Given.} #3\par
    \vspace{2pt}
    \textbf{Say the symbols as words.} #4\par
    \vspace{2pt}
    \textbf{For every item:} (1) say the whole statement in words; (2) lightly
    box its main relation, connective, or claim and mark its inputs or scope;
    (3) check whether it is complete and well formed here; (4) only then apply
    the definitions and circle an answer. \textbf{Decision key: N} = not a
    truth-evaluable proposition as written (incomplete, undefined, free
    variable, type error, or arity error); \textbf{T/F} = a complete proposition
    that is true/false. \textbf{N is not F. Do not circle before you parse.}}}
  \vspace{5pt}
  \begin{multicols}{2}
  \raggedcolumns
  #5
  \end{multicols}
}

\newcommand{\StageHeader}[2]{%
  \par\noindent
  {\color{CoreNavy}\footnotesize\bfseries #1}\par\nopagebreak
  {\color{CoreTeal}\scriptsize\itshape Reading lens: #2}\par\nopagebreak
  \vspace{1.5pt}
}

\newcommand{\StudentItem}[1]{%
  \stepcounter{tfitem}%
  \par\noindent
  \hangindent=3.25em\hangafter=1
  \makebox[1.65em][r]{\textbf{\thetfitem.}}\hspace{0.25em}%
  \makebox[1.25em][c]{\textbf{T}}\hspace{0.10em}%
  \makebox[1.25em][c]{\textbf{F}}\hspace{0.10em}%
  \makebox[1.25em][c]{\textbf{N}}\hspace{0.30em}%
  #1%
  \par\vspace{5.5pt}
}

\newcommand{\KeyItem}[5]{%
  \stepcounter{tfitem}%
  \par\noindent
  \hangindent=3.25em\hangafter=1
  \makebox[1.65em][r]{\textbf{\thetfitem.}}\hspace{0.25em}%
  \makebox[1.25em][c]{\textbf{T}}\hspace{0.10em}%
  \makebox[1.25em][c]{\textbf{F}}\hspace{0.10em}%
  \makebox[1.25em][c]{\textbf{N}}\hspace{0.30em}%
  #5%
  \par\nobreak
  {\scriptsize\color{CoreTeal}\hspace*{3.25em}%
    \textbf{[#1 --- #2]} #3%
    \if\relax\detokenize{#4}\relax\else\
      \textit{Minimal repair:} #4%
    \fi
  }\par
}

\begin{document}
\ifanswers\small\else\fontsize{10.8}{13.0}\selectfont\fi

\CoreForm{Logic}{A}{$p$ and $r$ are true; $q$ is false. Variables range over $\mathbb Z$. $P(n)$ means \emph{$n$ is even}, $Q(n)$ means \emph{$n>0$}, and $R(a,b)$ means \emph{$a<b$}. $P,Q$ are unary predicates and $R$ is binary.}{$\neg p$: not $p$; $p\land q$: $p$ and $q$; $p\lor q$: $p$ or $q$; $p\to q$: if $p$, then $q$; $p\leftrightarrow q$: $p$ exactly when $q$; $\forall$: for every; $\exists$: there exists; $\therefore$: therefore.}{
\StageHeader{1. Build and read propositions}{Locate every operator and all of its operands before assigning a truth value.}
\KeyItem{F}{F01: Formation}{The assigned truth value of $q$ is false.}{}{$q$}
\KeyItem{T}{F02: Formation}{Negation reverses the false value of $q$.}{}{$\neg q$}
\KeyItem{T}{F03: Formation}{Both conjuncts are true.}{}{$p\land r$}
\KeyItem{F}{F04: Formation}{Both disjuncts are false.}{}{$q\lor q$}
\KeyItem{F}{F05: Semantics}{A true antecedent and false consequent make it false.}{}{$p\to q$}
\StageHeader{2. Interpret connectives and equivalences}{Find the main connective first; then evaluate from the inside out.}
\KeyItem{T}{F06: Semantics}{The two sides have the same truth value.}{}{$p\leftrightarrow r$}
\KeyItem{N}{F07: Equivalence}{The negation sign has no operand.}{$p\lor\neg p$ is a tautology.}{$p\lor\neg$ is a tautology.}
\KeyItem{N}{F08: Equivalence}{The equivalence symbol is missing its right formula.}{$p\to q\equiv\neg p\lor q$.}{$p\to q\equiv$.}
\KeyItem{F}{F09: Equivalence}{The correct right side uses a conjunction.}{}{$\neg(p\lor q)\equiv\neg p\lor\neg q$.}
\KeyItem{T}{F10: Semantics}{Two independent truth values give $2^2=4$ rows.}{}{A truth table for $p,q$ has $4$ rows.}
\columnbreak
\StageHeader{3. Check predicates, types, and scope}{Check predicate arity, object types, and whether every variable is bound.}
\KeyItem{N}{F11: Arity/type}{A unary predicate symbol without an argument is not a proposition.}{$P(2)$.}{$P$.}
\KeyItem{N}{F12: Arity/type}{The binary predicate $R$ has only one argument.}{$R(1,2)$.}{$R(1)$.}
\KeyItem{F}{F13: Binding/scope}{Not every integer is even.}{}{$\forall n\,P(n)$.}
\KeyItem{T}{F14: Binding/scope}{Two is an even positive witness.}{}{$\exists n\,[P(n)\land Q(n)]$.}
\KeyItem{F}{F15: Binding/scope}{The quantifier on the right must be universal.}{}{$\neg\exists n\,P(n)\equiv\exists n\,\neg P(n)$.}
\StageHeader{4. Read quantified claims and arguments}{Read quantifiers left to right, mark their scope, and separate premises from conclusions.}
\KeyItem{N}{F16: Binding/scope}{The existential quantifier does not bind a variable.}{$\forall a\,\exists b\,R(a,b)$.}{$\forall a\,\exists\,R(a,b)$.}
\KeyItem{N}{F17: Proposition/term}{The open formula has a free variable and no fixed truth value.}{$P(2)$ or $\forall x\,P(x)$.}{$P(x)$.}
\KeyItem{F}{F18: Inference}{This is affirming the consequent.}{}{$r,\ p\to r\ \therefore p$ is valid.}
\KeyItem{N}{F19: Inference}{The disjunction is missing its right operand.}{$p\lor q,\ \neg p\ \therefore q$.}{$p\lor,\ \neg p\ \therefore q$ is valid.}
\KeyItem{T}{F20: Proposition/term}{This closed inequality is true.}{}{$1<2$.}
}
\CoreForm{Logic}{B}{$p$ and $r$ are true; $q$ is false. Variables range over $\mathbb Z$. $P(n)$ means \emph{$n$ is even}, $Q(n)$ means \emph{$n>0$}, and $R(a,b)$ means \emph{$a<b$}. $P,Q$ are unary predicates and $R$ is binary.}{$\neg p$: not $p$; $p\land q$: $p$ and $q$; $p\lor q$: $p$ or $q$; $p\to q$: if $p$, then $q$; $p\leftrightarrow q$: $p$ exactly when $q$; $\forall$: for every; $\exists$: there exists; $\therefore$: therefore.}{
\StageHeader{1. Build and read propositions}{Locate every operator and all of its operands before assigning a truth value.}
\KeyItem{N}{F01: Formation}{The binary connective is missing its right operand.}{$p\land q$}{$p\land$}
\KeyItem{F}{F02: Formation}{Negation reverses the true value of $p$.}{}{$\neg p$}
\KeyItem{F}{F03: Formation}{The second conjunct is false.}{}{$p\land q$}
\KeyItem{N}{F04: Formation}{The binary connective is missing its left operand.}{$p\lor q$}{$\lor q$}
\KeyItem{N}{F05: Semantics}{The conditional is missing its consequent.}{$p\to q$}{$p\to$}
\StageHeader{2. Interpret connectives and equivalences}{Find the main connective first; then evaluate from the inside out.}
\KeyItem{F}{F06: Semantics}{The two sides have different truth values.}{}{$p\leftrightarrow q$}
\KeyItem{T}{F07: Equivalence}{Excluded middle is true under every valuation.}{}{$p\lor\neg p$ is a tautology.}
\KeyItem{T}{F08: Equivalence}{This is the material-implication equivalence.}{}{$p\to q\equiv\neg p\lor q$.}
\KeyItem{N}{F09: Equivalence}{The conjunction inside the parentheses is incomplete.}{$\neg(p\land q)\equiv\neg p\lor\neg q$.}{$\neg(p\land)\equiv\neg p$.}
\KeyItem{F}{F10: Semantics}{Eight rows are needed for three atomic propositions, not two.}{}{A truth table for $p,q$ has $8$ rows.}
\columnbreak
\StageHeader{3. Check predicates, types, and scope}{Check predicate arity, object types, and whether every variable is bound.}
\KeyItem{T}{F11: Arity/type}{Two is even.}{}{$P(2)$.}
\KeyItem{T}{F12: Arity/type}{The inequality $1<2$ is true.}{}{$R(1,2)$.}
\KeyItem{N}{F13: Binding/scope}{The universal quantifier does not bind a variable.}{$\forall n\,P(n)$.}{$\forall\,P(n)$.}
\KeyItem{F}{F14: Binding/scope}{No integer can satisfy a predicate and its negation.}{}{$\exists n\,[P(n)\land\neg P(n)]$.}
\KeyItem{N}{F15: Binding/scope}{The quantifier does not bind a variable.}{$\neg\forall n\,P(n)$.}{$\neg\forall\,P(n)$.}
\StageHeader{4. Read quantified claims and arguments}{Read quantifiers left to right, mark their scope, and separate premises from conclusions.}
\KeyItem{T}{F16: Binding/scope}{For any integer $a$, choose $b=a+1$.}{}{$\forall a\,\exists b\,R(a,b)$.}
\KeyItem{T}{F17: Proposition/term}{All variables are bound and excluded middle applies.}{}{$\forall x\,[P(x)\lor\neg P(x)]$.}
\KeyItem{N}{F18: Inference}{The conditional premise is incomplete.}{$p,\ p\to r\ \therefore r$.}{$p,\ p\to\ \therefore r$ is valid.}
\KeyItem{T}{F19: Inference}{This is modus tollens.}{}{$p\to q,\ \neg q\ \therefore\neg p$ is valid.}
\KeyItem{F}{F20: Proposition/term}{This closed inequality is false.}{}{$2<1$.}
}
\CoreForm{Logic}{C}{$p$ and $r$ are true; $q$ is false. Variables range over $\mathbb Z$. $P(n)$ means \emph{$n$ is even}, $Q(n)$ means \emph{$n>0$}, and $R(a,b)$ means \emph{$a<b$}. $P,Q$ are unary predicates and $R$ is binary.}{$\neg p$: not $p$; $p\land q$: $p$ and $q$; $p\lor q$: $p$ or $q$; $p\to q$: if $p$, then $q$; $p\leftrightarrow q$: $p$ exactly when $q$; $\forall$: for every; $\exists$: there exists; $\therefore$: therefore.}{
\StageHeader{1. Build and read propositions}{Locate every operator and all of its operands before assigning a truth value.}
\KeyItem{T}{F01: Formation}{The assigned truth value of $p$ is true.}{}{$p$}
\KeyItem{N}{F02: Formation}{Negation has no operand.}{$\neg p$}{$\neg$}
\KeyItem{N}{F03: Formation}{Juxtaposing two proposition symbols is not a formation rule.}{$p\land q$}{$p\,q$}
\KeyItem{T}{F04: Formation}{At least one disjunct is true.}{}{$q\lor r$}
\KeyItem{T}{F05: Semantics}{A false antecedent makes the conditional true.}{}{$q\to p$}
\StageHeader{2. Interpret connectives and equivalences}{Find the main connective first; then evaluate from the inside out.}
\KeyItem{N}{F06: Semantics}{The conjunction has no left operand.}{$p\leftrightarrow(p\land q)$}{$p\leftrightarrow\land q$}
\KeyItem{F}{F07: Equivalence}{A proposition and its negation cannot both be true.}{}{$p\land\neg p$ is a tautology.}
\KeyItem{F}{F08: Equivalence}{The expressions differ when $p$ is false.}{}{$p\to q\equiv p\land q$.}
\KeyItem{T}{F09: Equivalence}{This is De Morgan's law for conjunction.}{}{$\neg(p\land q)\equiv\neg p\lor\neg q$.}
\KeyItem{N}{F10: Semantics}{The box is an unfilled exponent placeholder, not a number.}{A truth table for $p,q$ has $2^2$ rows.}{A truth table for $p,q$ has $2^{\Box}$ rows.}
\columnbreak
\StageHeader{3. Check predicates, types, and scope}{Check predicate arity, object types, and whether every variable is bound.}
\KeyItem{F}{F11: Arity/type}{Three is not even.}{}{$P(3)$.}
\KeyItem{F}{F12: Arity/type}{The inequality $2<1$ is false.}{}{$R(2,1)$.}
\KeyItem{T}{F13: Binding/scope}{For every integer, a proposition or its negation is true.}{}{$\forall n\,[P(n)\lor\neg P(n)]$.}
\KeyItem{N}{F14: Binding/scope}{The quantifier has no formula in its scope.}{$\exists n\,P(n)$.}{$\exists n$.}
\KeyItem{T}{F15: Binding/scope}{Negating a universal produces an existential with a negated predicate.}{}{$\neg\forall n\,P(n)\equiv\exists n\,\neg P(n)$.}
\StageHeader{4. Read quantified claims and arguments}{Read quantifiers left to right, mark their scope, and separate premises from conclusions.}
\KeyItem{F}{F16: Binding/scope}{No integer is greater than every integer.}{}{$\exists b\,\forall a\,R(a,b)$.}
\KeyItem{F}{F17: Proposition/term}{No integer makes a predicate and its negation both true.}{}{$\forall x\,[P(x)\land\neg P(x)]$.}
\KeyItem{T}{F18: Inference}{This is modus ponens.}{}{$p,\ p\to r\ \therefore r$ is valid.}
\KeyItem{F}{F19: Inference}{Disjunctive syllogism concludes $q$, not $\neg q$.}{}{$p\lor q,\ \neg p\ \therefore\neg q$ is valid.}
\KeyItem{N}{F20: Proposition/term}{This is a numerical term, not a proposition.}{$1+2=3$.}{$1+2$.}
}
\CoreForm{Arguments}{A}{$p,q,r$ denote propositions. A displayed argument lists premises before $\therefore$ and one conclusion after it. $\Gamma\models C$ means every valuation satisfying all premises in $\Gamma$ also satisfies $C$.}{$\neg$: not; $\land$: and; $\lor$: inclusive or; $p\to q$: if $p$, then $q$; $\Gamma\models C$: $C$ follows from premises $\Gamma$; $\therefore$: therefore.}{
\StageHeader{1. Recognize valid argument forms}{Separate the premises from the conclusion, then name the proposed inference pattern.}
\KeyItem{N}{F01: Valid form}{The argument has no conclusion.}{$p,\ p\to q\ \therefore q$.}{$p,\ p\to q\ \therefore$ is valid.}
\KeyItem{T}{F02: Valid form}{This is modus tollens.}{}{$p\to q,\ \neg q\ \therefore\neg p$ is valid.}
\KeyItem{N}{F03: Valid form}{The disjunction is missing an operand.}{$p\lor q,\ \neg p\ \therefore q$.}{$p\lor,\ \neg p\ \therefore q$ is valid.}
\KeyItem{F}{F04: Valid form}{The shared consequent does not link $p$ to $r$.}{}{$p\to q,\ r\to q\ \therefore p\to r$ is valid.}
\KeyItem{T}{F05: Valid form}{Conjunction combines two established propositions.}{}{$p,\ q\ \therefore p\land q$ is valid.}
\StageHeader{2. Separate validity, soundness, and consistency}{Ask whether the premises could all be true while the conclusion is false.}
\KeyItem{F}{F06: Valid form}{A disjunction does not identify which disjunct is true.}{}{$p\lor q\ \therefore p$ is valid.}
\KeyItem{N}{F07: Validity/soundness}{The conjunction has no right-hand condition.}{Soundness $=$ valid $\land$ true premises.}{Soundness $=$ valid $\land$.}
\KeyItem{F}{F08: Validity/soundness}{That valuation does not demonstrate invalidity.}{}{A countermodel makes one premise false and the conclusion true.}
\KeyItem{T}{F09: Validity/soundness}{Validity concerns the inferential connection, not conclusion truth alone.}{}{An invalid argument can have a true conclusion.}
\KeyItem{N}{F10: Consistency}{The second expression is not a complete proposition.}{$p$ and $\neg p$ are inconsistent.}{$p$ and $\neg$ are inconsistent.}
\columnbreak
\StageHeader{3. Read assumptions, flaws, and conditionals}{Identify exactly what the conclusion needs and what each premise actually supplies.}
\KeyItem{F}{F11: Assumptions}{Relevance to the inferential link is essential.}{}{A necessary assumption is merely an unrelated true statement.}
\KeyItem{N}{F12: Flaws}{The displayed argument is missing its conclusion.}{$p\therefore p$.}{$p\therefore$ because $p$.}
\KeyItem{T}{F13: Causal reasoning}{Alternative causes and reverse direction remain possible.}{}{A correlation alone does not establish which variable causes the other.}
\KeyItem{N}{F14: Generalization}{The subset relation is missing the population on the right.}{sample $\subseteq$ population.}{sample $\subseteq$.}
\KeyItem{F}{F15: Conditionals}{The conditional does not say that $q$ guarantees $p$.}{}{In $p\to q$, $q$ is sufficient for $p$.}
\StageHeader{4. Test what the information supports}{Use only the stated information; distinguish support from proof and relevance from truth.}
\KeyItem{F}{F16: Strengthen}{Additional support need not amount to a proof.}{}{A strengthening answer must independently prove the conclusion.}
\KeyItem{N}{F17: Weaken}{The disjunction has no second alternative.}{alternative A $\lor$ alternative B.}{alternative $\lor$.}
\KeyItem{T}{F18: Must be true}{It must hold in every case allowed by those facts.}{}{A must-be-true answer follows from all the supplied facts.}
\KeyItem{T}{F19: Comparison}{That is a genuine disagreement.}{}{A point-at-issue answer states something one speaker accepts and the other rejects.}
\KeyItem{F}{F20: Argument structure}{Matching topic words does not determine inferential structure.}{}{Parallel reasoning requires the same nouns in both arguments.}
}
\CoreForm{Arguments}{B}{$p,q,r$ denote propositions. A displayed argument lists premises before $\therefore$ and one conclusion after it. $\Gamma\models C$ means every valuation satisfying all premises in $\Gamma$ also satisfies $C$.}{$\neg$: not; $\land$: and; $\lor$: inclusive or; $p\to q$: if $p$, then $q$; $\Gamma\models C$: $C$ follows from premises $\Gamma$; $\therefore$: therefore.}{
\StageHeader{1. Recognize valid argument forms}{Separate the premises from the conclusion, then name the proposed inference pattern.}
\KeyItem{T}{F01: Valid form}{This is modus ponens.}{}{$p,\ p\to q\ \therefore q$ is valid.}
\KeyItem{F}{F02: Valid form}{This is denying the antecedent.}{}{$p\to q,\ \neg p\ \therefore\neg q$ is valid.}
\KeyItem{T}{F03: Valid form}{This is disjunctive syllogism.}{}{$p\lor q,\ \neg p\ \therefore q$ is valid.}
\KeyItem{N}{F04: Valid form}{The second premise is incomplete.}{$p\to q,\ q\to r\ \therefore p\to r$.}{$p\to q,\ q\to\ \therefore p\to r$ is valid.}
\KeyItem{F}{F05: Valid form}{The premise provides no support for $q$.}{}{$p\ \therefore p\land q$ is valid.}
\StageHeader{2. Separate validity, soundness, and consistency}{Ask whether the premises could all be true while the conclusion is false.}
\KeyItem{N}{F06: Valid form}{The conclusion's disjunction has no left operand.}{$p\ \therefore p\lor q$.}{$p\ \therefore\lor q$ is valid.}
\KeyItem{T}{F07: Validity/soundness}{That is the definition of soundness.}{}{A sound argument is valid and has true premises.}
\KeyItem{N}{F08: Validity/soundness}{Semantic entailment is missing its conclusion.}{$\Gamma\models C$.}{$\Gamma\models$.}
\KeyItem{F}{F09: Validity/soundness}{A conclusion may be accidentally true.}{}{A true conclusion makes every argument for it valid.}
\KeyItem{T}{F10: Consistency}{No valuation makes both true.}{}{$p$ and $\neg p$ are inconsistent.}
\columnbreak
\StageHeader{3. Read assumptions, flaws, and conditionals}{Identify exactly what the conclusion needs and what each premise actually supplies.}
\KeyItem{N}{F11: Assumptions}{The because marker is missing the supporting proposition.}{$q$ because $p$.}{premise $\because$.}
\KeyItem{T}{F12: Flaws}{That is the defining defect of circular reasoning.}{}{A circular argument uses its conclusion as support for itself.}
\KeyItem{F}{F13: Causal reasoning}{Association is not sufficient evidence of mechanism.}{}{A correlation by itself proves a causal mechanism.}
\KeyItem{T}{F14: Generalization}{Representativeness affects how well the sample supports the claim.}{}{A representative sample is relevant to a population generalization.}
\KeyItem{N}{F15: Conditionals}{The conditional has no consequent.}{In $p\to q$, $q$ is necessary for $p$.}{In $p\to$, the consequent is necessary for $p$.}
\StageHeader{4. Test what the information supports}{Use only the stated information; distinguish support from proof and relevance from truth.}
\KeyItem{N}{F16: Strengthen}{The support marker is missing a supporting claim.}{premises support conclusion.}{premises $\uparrow$ conclusion $\because$.}
\KeyItem{T}{F17: Weaken}{An alternative explanation can reduce causal support.}{}{A weakening answer may expose an alternative explanation.}
\KeyItem{F}{F18: Must be true}{Compatibility is weaker than entailment.}{}{A statement that is merely compatible with the facts must be true.}
\KeyItem{F}{F19: Comparison}{They may use different examples while accepting the same claim.}{}{Different examples guarantee that two speakers disagree.}
\KeyItem{N}{F20: Argument structure}{Arithmetic addition is not defined for proposition symbols here.}{$p\land q\ \therefore r$.}{$p+q\ \therefore r$.}
}
\CoreForm{Arguments}{C}{$p,q,r$ denote propositions. A displayed argument lists premises before $\therefore$ and one conclusion after it. $\Gamma\models C$ means every valuation satisfying all premises in $\Gamma$ also satisfies $C$.}{$\neg$: not; $\land$: and; $\lor$: inclusive or; $p\to q$: if $p$, then $q$; $\Gamma\models C$: $C$ follows from premises $\Gamma$; $\therefore$: therefore.}{
\StageHeader{1. Recognize valid argument forms}{Separate the premises from the conclusion, then name the proposed inference pattern.}
\KeyItem{F}{F01: Valid form}{This is affirming the consequent.}{}{$q,\ p\to q\ \therefore p$ is valid.}
\KeyItem{N}{F02: Valid form}{The conditional premise is incomplete.}{$p\to q,\ \neg q\ \therefore\neg p$.}{$p\to,\ \neg q\ \therefore\neg p$ is valid.}
\KeyItem{F}{F03: Valid form}{An inclusive disjunction permits both disjuncts to be true.}{}{$p\lor q,\ p\ \therefore\neg q$ is valid.}
\KeyItem{T}{F04: Valid form}{This is hypothetical syllogism.}{}{$p\to q,\ q\to r\ \therefore p\to r$ is valid.}
\KeyItem{N}{F05: Valid form}{The conclusion is missing a conjunct.}{$p,\ q\ \therefore p\land q$.}{$p,\ q\ \therefore p\land$ is valid.}
\StageHeader{2. Separate validity, soundness, and consistency}{Ask whether the premises could all be true while the conclusion is false.}
\KeyItem{T}{F06: Valid form}{This is addition.}{}{$p\ \therefore p\lor q$ is valid.}
\KeyItem{F}{F07: Validity/soundness}{Validity concerns truth preservation, not the actual truth of premises.}{}{A valid argument must have true premises.}
\KeyItem{T}{F08: Validity/soundness}{Such a valuation directly violates truth preservation.}{}{A countermodel to validity makes every premise true and the conclusion false.}
\KeyItem{N}{F09: Validity/soundness}{The conclusion after the therefore sign is missing.}{$p\ \therefore p$.}{true $\therefore$.}
\KeyItem{F}{F10: Consistency}{Any valuation with $p$ true satisfies both.}{}{$p$ and $p\lor q$ are inconsistent.}
\columnbreak
\StageHeader{3. Read assumptions, flaws, and conditionals}{Identify exactly what the conclusion needs and what each premise actually supplies.}
\KeyItem{T}{F11: Assumptions}{Removing a necessary assumption breaks the inferential link.}{}{A necessary assumption must hold for the stated support to work.}
\KeyItem{F}{F12: Flaws}{Truth of the conclusion does not remove circular support.}{}{An argument becomes noncircular whenever its conclusion is true.}
\KeyItem{N}{F13: Causal reasoning}{The function-like notation has an unmatched delimiter and missing argument.}{$\operatorname{cause}(p,q)$.}{cause$(p,q,$.}
\KeyItem{F}{F14: Generalization}{Both can materially affect its strength.}{}{Sample size and selection method never affect a generalization.}
\KeyItem{T}{F15: Conditionals}{If $p$ occurs, $q$ must occur.}{}{In $p\to q$, $q$ is necessary for $p$.}
\StageHeader{4. Test what the information supports}{Use only the stated information; distinguish support from proof and relevance from truth.}
\KeyItem{T}{F16: Strengthen}{That is the task of a strengthen question.}{}{A strengthening answer adds support to the conclusion relative to the premises.}
\KeyItem{F}{F17: Weaken}{Reducing support is enough; refutation is not required.}{}{A weakening answer must establish the exact negation of the conclusion.}
\KeyItem{N}{F18: Must be true}{The entailment symbol is missing the proposed conclusion.}{facts $\models C$.}{facts $\models$.}
\KeyItem{N}{F19: Comparison}{The biconditional is missing its right-hand statement.}{speaker A accepts $p$; speaker B rejects $p$.}{speaker A $\leftrightarrow$.}
\KeyItem{T}{F20: Argument structure}{Logical form abstracts away from subject matter.}{}{Parallel arguments can share a logical form while discussing different topics.}
}
\CoreForm{Sets}{A}{$U=\{1,2,3,4\}$, $A=\{1,2\}$, $B=\{2,3\}$, and $C=A$. Complements are relative to $U$. Let $f:A\to U$ satisfy $f(1)=2$, $f(2)=4$, and let $R=\{(1,1),(2,2)\}$ be a relation on $A$.}{$x\in A$: $x$ is an element of $A$; $A\subseteq B$: every $A$-element is in $B$; $\cup$: in either set; $\cap$: in both; $A\setminus B$: in $A$ but not $B$; $A^c$: outside $A$; $\mathcal P(A)$: the set of all subsets of $A$.}{
\StageHeader{1. Read membership, inclusion, and size}{Name the object on each side and distinguish membership from subset inclusion.}
\KeyItem{F}{F01: Membership}{Three is not an element of $A$.}{}{$3\in A$.}
\KeyItem{N}{F02: Subset}{The subset relation is missing its right operand.}{$A\subseteq U$.}{$A\subseteq$.}
\KeyItem{N}{F03: Equality}{The set delimiter is not closed.}{$A=\{1,2\}$.}{$A=\{1,2$.}
\KeyItem{F}{F04: Empty set}{Neither element of $A$ is the empty set.}{}{$\varnothing\in A$.}
\KeyItem{F}{F05: Cardinality}{The set $A$ does not have three elements.}{}{$|A|=3$.}
\StageHeader{2. Interpret set operations}{Translate each operation into an element-by-element condition before computing.}
\KeyItem{T}{F06: Union}{The union contains everything in either set.}{}{$A\cup B=\{1,2,3\}$.}
\KeyItem{N}{F07: Intersection}{Intersection is missing a set and membership is missing an element.}{$A\cap B=\{2\}$.}{$A\cap\in B$.}
\KeyItem{F}{F08: Difference}{Two lies in both sets and is removed.}{}{$A\setminus B=\{2\}$.}
\KeyItem{T}{F09: Complement}{These are the elements of $U$ outside $A$.}{}{$A^c=\{3,4\}$.}
\KeyItem{T}{F10: Set laws}{This is De Morgan's law for sets.}{}{$(A\cup B)^c=A^c\cap B^c$.}
\columnbreak
\StageHeader{3. Read products and relations}{Check the type of every object: element, ordered pair, set, or relation.}
\KeyItem{N}{F11: Power set}{The power-set operator has no set argument.}{$\mathcal P(A)$.}{$\mathcal P()=A$.}
\KeyItem{T}{F12: Power set}{The empty set is a subset of $A$.}{}{$\varnothing\in\mathcal P(A)$.}
\KeyItem{F}{F13: Ordered pairs}{The first coordinates differ.}{}{$(1,2)=(2,1)$.}
\KeyItem{N}{F14: Cartesian product}{The product is missing a set and membership is missing an element.}{$A\times B$.}{$A\times\in B$.}
\KeyItem{T}{F15: Relations}{Both $(1,1)$ and $(2,2)$ belong to $R$.}{}{$R$ is reflexive on $A$.}
\StageHeader{4. Read functions, variables, and types}{Bind every variable and distinguish a function value from a claim about the function.}
\KeyItem{N}{F16: Relations}{The second ordered pair is incomplete.}{$(1,2)\in R\Rightarrow(2,1)\in R$.}{$(1,2)\in R\Rightarrow(2,)\in R$.}
\KeyItem{F}{F17: Functions}{The context specifies $f(2)=4$.}{}{$f(2)=2$.}
\KeyItem{F}{F18: Functions}{The values $1$ and $3$ are not outputs.}{}{$f$ is surjective onto $U$.}
\KeyItem{T}{F19: Binding/type}{Every element of $A$ lies in $U$.}{}{$\forall x\in A,\ x\in U$.}
\KeyItem{N}{F20: Proposition/term}{This is a set-valued term, not a proposition.}{$A\cap B=\{2\}$.}{$A\cap B$.}
}
\CoreForm{Sets}{B}{$U=\{1,2,3,4\}$, $A=\{1,2\}$, $B=\{2,3\}$, and $C=A$. Complements are relative to $U$. Let $f:A\to U$ satisfy $f(1)=2$, $f(2)=4$, and let $R=\{(1,1),(2,2)\}$ be a relation on $A$.}{$x\in A$: $x$ is an element of $A$; $A\subseteq B$: every $A$-element is in $B$; $\cup$: in either set; $\cap$: in both; $A\setminus B$: in $A$ but not $B$; $A^c$: outside $A$; $\mathcal P(A)$: the set of all subsets of $A$.}{
\StageHeader{1. Read membership, inclusion, and size}{Name the object on each side and distinguish membership from subset inclusion.}
\KeyItem{N}{F01: Membership}{Membership is missing the set on the right.}{$1\in A$.}{$1\in$.}
\KeyItem{T}{F02: Subset}{Every element of $A$ lies in $U$.}{}{$A\subseteq U$.}
\KeyItem{T}{F03: Equality}{The context defines $C$ to equal $A$.}{}{$A=C$.}
\KeyItem{N}{F04: Empty set}{The subset relation has no right operand.}{$\varnothing\subseteq A$.}{$\varnothing\subseteq$.}
\KeyItem{N}{F05: Cardinality}{The equality has no right-hand value.}{$|A|=2$.}{$|A|=$.}
\StageHeader{2. Interpret set operations}{Translate each operation into an element-by-element condition before computing.}
\KeyItem{F}{F06: Union}{The display gives the intersection, not the union.}{}{$A\cup B=\{2\}$.}
\KeyItem{T}{F07: Intersection}{Two is the only shared element.}{}{$A\cap B=\{2\}$.}
\KeyItem{N}{F08: Difference}{Set difference is missing its second set.}{$A\setminus B$.}{$A\setminus$.}
\KeyItem{F}{F09: Complement}{That set is $A$, not its complement.}{}{$A^c=\{1,2\}$.}
\KeyItem{F}{F10: Set laws}{The operation on the right must switch to intersection.}{}{$(A\cup B)^c=A^c\cup B^c$.}
\columnbreak
\StageHeader{3. Read products and relations}{Check the type of every object: element, ordered pair, set, or relation.}
\KeyItem{T}{F11: Power set}{A two-element set has $2^2=4$ subsets.}{}{$|\mathcal P(A)|=4$.}
\KeyItem{F}{F12: Power set}{The set $\{3\}$ is not a subset of $A$.}{}{$\{3\}\in\mathcal P(A)$.}
\KeyItem{N}{F13: Ordered pairs}{The ordered pair is missing its second coordinate.}{$(1,2)\in A\times B$.}{$(1,)\in A\times B$.}
\KeyItem{T}{F14: Cartesian product}{There are $2\cdot2=4$ ordered pairs.}{}{$|A\times B|=4$.}
\KeyItem{F}{F15: Relations}{Every element is related to itself.}{}{$R$ is irreflexive on $A$.}
\StageHeader{4. Read functions, variables, and types}{Bind every variable and distinguish a function value from a claim about the function.}
\KeyItem{T}{F16: Relations}{Every pair in $R$ is diagonal and reverses to itself.}{}{$R$ is symmetric on $A$.}
\KeyItem{N}{F17: Functions}{The function $f$ has one input, not two.}{$f(2)=4$.}{$f(1,2)=4$.}
\KeyItem{N}{F18: Functions}{The conjunction is incomplete, and $f(A)$ is not defined by the stated elementwise notation.}{$f(1)=2$.}{$f(A)=2\land$.}
\KeyItem{F}{F19: Binding/type}{Three and four are counterexamples.}{}{$\forall x\in U,\ x\in A$.}
\KeyItem{T}{F20: Proposition/term}{Two is in both sets.}{}{$2\in A\cap B$.}
}
\CoreForm{Sets}{C}{$U=\{1,2,3,4\}$, $A=\{1,2\}$, $B=\{2,3\}$, and $C=A$. Complements are relative to $U$. Let $f:A\to U$ satisfy $f(1)=2$, $f(2)=4$, and let $R=\{(1,1),(2,2)\}$ be a relation on $A$.}{$x\in A$: $x$ is an element of $A$; $A\subseteq B$: every $A$-element is in $B$; $\cup$: in either set; $\cap$: in both; $A\setminus B$: in $A$ but not $B$; $A^c$: outside $A$; $\mathcal P(A)$: the set of all subsets of $A$.}{
\StageHeader{1. Read membership, inclusion, and size}{Name the object on each side and distinguish membership from subset inclusion.}
\KeyItem{T}{F01: Membership}{One is listed as an element of $A$.}{}{$1\in A$.}
\KeyItem{F}{F02: Subset}{Three belongs to $B$ but not to $A$.}{}{$B\subseteq A$.}
\KeyItem{F}{F03: Equality}{The sets have different elements.}{}{$A=B$.}
\KeyItem{T}{F04: Empty set}{The empty set is a subset of every set.}{}{$\varnothing\subseteq A$.}
\KeyItem{T}{F05: Cardinality}{The set $A$ has two elements.}{}{$|A|=2$.}
\StageHeader{2. Interpret set operations}{Translate each operation into an element-by-element condition before computing.}
\KeyItem{N}{F06: Union}{Union is missing its second set.}{$A\cup B=\{1,2,3\}$.}{$A\cup=\{1,2,3\}$.}
\KeyItem{F}{F07: Intersection}{The display gives the union.}{}{$A\cap B=\{1,2,3\}$.}
\KeyItem{T}{F08: Difference}{One is in $A$ but not $B$.}{}{$A\setminus B=\{1\}$.}
\KeyItem{N}{F09: Complement}{The union inside the parentheses is incomplete.}{$(A\cup B)^c$.}{$(A\cup)^c$.}
\KeyItem{N}{F10: Set laws}{The intersection is missing its second set.}{$(A\cup B)^c=A^c\cap B^c$.}{$(A\cup B)^c=A^c\cap$.}
\columnbreak
\StageHeader{3. Read products and relations}{Check the type of every object: element, ordered pair, set, or relation.}
\KeyItem{F}{F11: Power set}{The power set contains more than the singleton subsets.}{}{$|\mathcal P(A)|=2$.}
\KeyItem{N}{F12: Power set}{Membership is missing the set on the right.}{$\varnothing\in\mathcal P(A)$.}{$\mathcal P(A)\in$.}
\KeyItem{T}{F13: Ordered pairs}{Ordered pairs are equal only when corresponding entries agree.}{}{$(1,2)\ne(2,1)$.}
\KeyItem{F}{F14: Cartesian product}{The cardinalities multiply.}{}{$|A\times B|=2$.}
\KeyItem{N}{F15: Relations}{The binary relation notation is missing a second element.}{$1R1$.}{$1R$.}
\StageHeader{4. Read functions, variables, and types}{Bind every variable and distinguish a function value from a claim about the function.}
\KeyItem{F}{F16: Relations}{Its reverse pair $(2,1)$ is missing.}{}{$\{(1,2)\}$ is symmetric on $A$.}
\KeyItem{T}{F17: Functions}{This value is specified in the context.}{}{$f(1)=2$.}
\KeyItem{T}{F18: Functions}{Its two distinct inputs have distinct outputs.}{}{$f$ is injective.}
\KeyItem{N}{F19: Binding/type}{The variable $x$ is neither assigned nor bound.}{$\forall x\in A,\ x\in U$.}{$x\in A$.}
\KeyItem{F}{F20: Proposition/term}{One is not in $B$.}{}{$1\in A\cap B$.}
}
\CoreForm{Proofs}{A}{Variables range over $\mathbb Z$ when bound. $P$ and $Q$ denote propositions, and $P(n)$ denotes a proposition about an integer $n$. A formula with an unbound variable is not a complete proposition.}{$\forall$: for every; $\exists$: there exists; $P\to Q$: if $P$, then $Q$; $P\leftrightarrow Q$: both directions; a \emph{witness} proves an existential; a \emph{counterexample} disproves a universal.}{
\StageHeader{1. Identify the proof strategy}{State what may be assumed and what must be shown for the named strategy.}
\KeyItem{T}{F01: Direct proof}{That is the standard direct-proof structure.}{}{A direct proof of $P\to Q$ may assume $P$ and derive $Q$.}
\KeyItem{N}{F02: Contrapositive}{The conditional's consequent is incomplete.}{$\neg Q\to\neg P$.}{$\neg Q\to\neg$.}
\KeyItem{T}{F03: Contradiction}{A contradiction rules out the negation of the goal.}{}{To prove $P$ by contradiction, one may assume $\neg P$ and derive $\bot$.}
\KeyItem{F}{F04: Cases}{One case does not establish a general result.}{}{A proof by cases is complete after one convenient case.}
\KeyItem{N}{F05: Biconditional}{The biconditional is missing its right formula.}{$P\leftrightarrow Q$.}{$P\leftrightarrow$.}
\StageHeader{2. Read quantifiers, witnesses, and counterexamples}{Match universal claims with arbitrary objects and existential claims with witnesses.}
\KeyItem{F}{F06: Universal proof}{That proves at most a special case.}{}{A proof of a universal statement may choose only an object with a convenient extra property.}
\KeyItem{F}{F07: Existential proof}{That is stronger than an existential claim.}{}{To prove $\exists n\,P(n)$, one must prove $P(n)$ for every integer.}
\KeyItem{N}{F08: Instantiation}{The universal premise has no formula in its scope.}{From $\forall n\,P(n)$, infer $P(7)$.}{From $\forall n$, infer $P(7)$.}
\KeyItem{T}{F09: Instantiation}{The fresh symbol denotes some unspecified witness.}{}{From $\exists n\,P(n)$, one may introduce a fresh symbol $c$ with $P(c)$.}
\KeyItem{F}{F10: Counterexample}{A confirming example does not refute the universal.}{}{One integer $c$ with $P(c)$ disproves $\forall n\,P(n)$.}
\columnbreak
\StageHeader{3. Expand definitions inside proofs}{Expand the relevant definition before deciding whether a displayed line proves anything.}
\KeyItem{F}{F11: Definitions}{The displayed equality is false and gives no witness.}{}{$9=2(4)$ proves that $9$ is even.}
\KeyItem{N}{F12: Definitions}{This is an integer-valued term, not a proposition.}{$9=2(4)+1$.}{$2k+1$.}
\KeyItem{T}{F13: Definitions}{It is a quotient of integers with nonzero denominator.}{}{$\frac{3}{4}$ is rational.}
\KeyItem{F}{F14: Definitions}{No integer multiplied by six equals twenty.}{}{$6\mid20$.}
\KeyItem{N}{F15: Parity}{The equality has no parity class on the right.}{even $+$ odd $=$ odd.}{even $+$ odd $=$.}
\StageHeader{4. Audit proof language and logical structure}{Separate established facts from goals, conjectures, and incomplete proof language.}
\KeyItem{T}{F16: Parity}{$(2j+1)(2k+1)=2(2jk+j+k)+1$.}{}{The product of two odd integers is odd.}
\KeyItem{F}{F17: Proof language}{A conjecture is a proposed statement not yet proved.}{}{A conjecture is already established by proof.}
\KeyItem{T}{F18: Conditionals}{If $P$ holds, $Q$ must follow.}{}{In $P\to Q$, $Q$ is necessary for $P$.}
\KeyItem{N}{F19: Proof validity}{The therefore marker has no conclusion.}{$P\because P\therefore P$.}{$P\because P\therefore$.}
\KeyItem{N}{F20: Binding/type}{The variable $n$ is free and has not been assigned.}{$\forall n\in\mathbb Z,\ n^2\ge0$.}{$n^2\ge0$.}
}
\CoreForm{Proofs}{B}{Variables range over $\mathbb Z$ when bound. $P$ and $Q$ denote propositions, and $P(n)$ denotes a proposition about an integer $n$. A formula with an unbound variable is not a complete proposition.}{$\forall$: for every; $\exists$: there exists; $P\to Q$: if $P$, then $Q$; $P\leftrightarrow Q$: both directions; a \emph{witness} proves an existential; a \emph{counterexample} disproves a universal.}{
\StageHeader{1. Identify the proof strategy}{State what may be assumed and what must be shown for the named strategy.}
\KeyItem{F}{F01: Direct proof}{That proves the converse, not the stated conditional.}{}{A direct proof of $P\to Q$ may assume $Q$ and derive $P$.}
\KeyItem{T}{F02: Contrapositive}{A conditional is logically equivalent to its contrapositive.}{}{Proving $\neg Q\to\neg P$ proves $P\to Q$.}
\KeyItem{F}{F03: Contradiction}{Restating the goal does not derive a contradiction.}{}{To prove $P$ by contradiction, assume $P$ and stop.}
\KeyItem{N}{F04: Cases}{The disjunction of cases is incomplete.}{case 1 $\lor$ case 2.}{case 1 $\lor$.}
\KeyItem{T}{F05: Biconditional}{The two directions establish equivalence.}{}{To prove $P\leftrightarrow Q$, it is enough to prove $P\to Q$ and $Q\to P$.}
\StageHeader{2. Read quantifiers, witnesses, and counterexamples}{Match universal claims with arbitrary objects and existential claims with witnesses.}
\KeyItem{N}{F06: Universal proof}{The equality is incomplete.}{Let $n\in\mathbb Z$ be arbitrary.}{Let $n=$ be arbitrary.}
\KeyItem{N}{F07: Existential proof}{The existential quantifier does not bind a variable.}{$\exists n\,P(n)$.}{$\exists\,P(n)$.}
\KeyItem{T}{F08: Instantiation}{Universal instantiation permits any object in the domain.}{}{From $\forall n\,P(n)$, one may infer $P(7)$.}
\KeyItem{F}{F09: Instantiation}{An existential witness does not represent all objects.}{}{From $\exists n\,P(n)$, one may treat the witness as every integer.}
\KeyItem{N}{F10: Counterexample}{The universal claim being refuted is missing.}{$c$ is a counterexample to $\forall n\,P(n)$.}{$c$ is a counterexample to.}
\columnbreak
\StageHeader{3. Expand definitions inside proofs}{Expand the relevant definition before deciding whether a displayed line proves anything.}
\KeyItem{N}{F11: Definitions}{This is an integer-valued term, not a proposition.}{$6=2(3)$.}{$2k$.}
\KeyItem{T}{F12: Definitions}{It has the required form with an integer witness.}{}{$9=2(4)+1$ proves that $9$ is odd.}
\KeyItem{F}{F13: Definitions}{The square root of two is irrational.}{}{$\sqrt2$ is rational.}
\KeyItem{N}{F14: Definitions}{Divisibility is missing its right operand.}{$4\mid20$.}{$4\mid$.}
\KeyItem{T}{F15: Parity}{If $a=2j$ and $b=2k$, then $a+b=2(j+k)$.}{}{The sum of two even integers is even.}
\StageHeader{4. Audit proof language and logical structure}{Separate established facts from goals, conjectures, and incomplete proof language.}
\KeyItem{F}{F16: Parity}{The even factor makes the product even.}{}{The product of an even integer and an odd integer is odd.}
\KeyItem{N}{F17: Proof language}{The implication has no conclusion.}{theorem $\Rightarrow$ proved statement.}{theorem $\Rightarrow$.}
\KeyItem{F}{F18: Conditionals}{The conditional does not allow the converse inference.}{}{In $P\to Q$, $Q$ is sufficient for $P$.}
\KeyItem{T}{F19: Proof validity}{The conclusion cannot serve as its own independent support.}{}{Using the desired conclusion as an unproved premise is circular.}
\KeyItem{T}{F20: Binding/type}{Every integer square is nonnegative.}{}{$\forall n\in\mathbb Z,\ n^2\ge0$.}
}
\CoreForm{Proofs}{C}{Variables range over $\mathbb Z$ when bound. $P$ and $Q$ denote propositions, and $P(n)$ denotes a proposition about an integer $n$. A formula with an unbound variable is not a complete proposition.}{$\forall$: for every; $\exists$: there exists; $P\to Q$: if $P$, then $Q$; $P\leftrightarrow Q$: both directions; a \emph{witness} proves an existential; a \emph{counterexample} disproves a universal.}{
\StageHeader{1. Identify the proof strategy}{State what may be assumed and what must be shown for the named strategy.}
\KeyItem{N}{F01: Direct proof}{The proof fragment has no derived conclusion.}{Assume $P$; therefore $Q$.}{Assume $P$; therefore.}
\KeyItem{F}{F02: Contrapositive}{The converse need not be equivalent to the original.}{}{Proving $Q\to P$ proves $P\to Q$.}
\KeyItem{N}{F03: Contradiction}{The implication has no conclusion.}{$\neg P\Rightarrow\bot$.}{$\neg P\Rightarrow$.}
\KeyItem{T}{F04: Cases}{Otherwise the uncovered cases remain unproved.}{}{A proof by cases must cover every possibility allowed by the hypotheses.}
\KeyItem{F}{F05: Biconditional}{One direction does not establish a biconditional.}{}{To prove $P\leftrightarrow Q$, it is enough to prove only $P\to Q$.}
\StageHeader{2. Read quantifiers, witnesses, and counterexamples}{Match universal claims with arbitrary objects and existential claims with witnesses.}
\KeyItem{T}{F06: Universal proof}{The argument must not depend on a special choice.}{}{A proof of a universal statement may begin with an arbitrary object in the domain.}
\KeyItem{T}{F07: Existential proof}{An existential claim asks for at least one example.}{}{To prove $\exists n\,P(n)$, one valid witness is sufficient.}
\KeyItem{F}{F08: Instantiation}{The existential witness need not be the preferred object.}{}{From $\exists n\,P(n)$, one may infer $P(7)$ without further information.}
\KeyItem{N}{F09: Instantiation}{The existential statement is missing its predicate.}{Let $c$ witness $\exists n\,P(n)$.}{Let $c$ witness $\exists n$.}
\KeyItem{T}{F10: Counterexample}{A universal claim fails at any counterexample.}{}{One integer $c$ with $\neg P(c)$ disproves $\forall n\,P(n)$.}
\columnbreak
\StageHeader{3. Expand definitions inside proofs}{Expand the relevant definition before deciding whether a displayed line proves anything.}
\KeyItem{T}{F11: Definitions}{It exhibits the required integer multiplier.}{}{$6=2(3)$ proves that $6$ is even.}
\KeyItem{F}{F12: Definitions}{The displayed equality is false.}{}{$6=2(3)+1$ proves that $6$ is odd.}
\KeyItem{N}{F13: Definitions}{Division by zero is undefined, so the displayed term is not defined.}{$\frac{3}{4}$ is rational.}{$\frac{3}{0}$ is rational.}
\KeyItem{T}{F14: Definitions}{Twenty equals $4(5)$.}{}{$4\mid20$.}
\KeyItem{F}{F15: Parity}{If $a=2j+1$ and $b=2k+1$, then $a+b=2(j+k+1)$ is even.}{}{The sum of two odd integers is odd.}
\StageHeader{4. Audit proof language and logical structure}{Separate established facts from goals, conjectures, and incomplete proof language.}
\KeyItem{N}{F16: Parity}{Multiplication is missing its second factor.}{odd $\cdot$ odd $=$ odd.}{odd $\cdot$.}
\KeyItem{T}{F17: Proof language}{That is the role of a theorem.}{}{A theorem is a statement established by proof.}
\KeyItem{N}{F18: Conditionals}{The conditional has no consequent.}{In $P\to Q$, $Q$ is necessary for $P$.}{In $P\to$, the consequent is necessary.}
\KeyItem{F}{F19: Proof validity}{Repetition without justification is not a derivation.}{}{A proof is complete when its final claim merely repeats an unsupported first claim.}
\KeyItem{F}{F20: Binding/type}{No integer square is negative.}{}{$\forall n\in\mathbb Z,\ n^2<0$.}
}
\CoreForm{Divisibility}{A}{Let $a,b,c$ be arbitrary integers and let $m\ge2$ be a fixed integer. All other variables are integers when quantified or assigned. In $a\mid b$, the left integer divides the right integer. In $a\equiv_m b$, the modulus satisfies $m\ge2$. A base-$b$ numeral may use only digits $0,\ldots,b-1$.}{$a\mid b$: $a$ divides $b$; $a\nmid b$: $a$ does not divide $b$; $a\equiv b\pmod m$: $a$ and $b$ have the same remainder modulo $m$; $\gcd(a,b)$: greatest common divisor; $(d_k\cdots d_0)_b$: a base-$b$ numeral.}{
\StageHeader{1. Read divisibility from its definition}{Read $a\mid b$ as an existential equation: $b=ak$ for some integer $k$.}
\KeyItem{T}{F01: Divisibility}{Twelve equals $3(4)$.}{}{$3\mid12$.}
\KeyItem{T}{F02: Definition}{This is the definition of divisibility.}{}{$4\mid20\iff\exists k\in\mathbb Z,\ 20=4k$.}
\KeyItem{N}{F03: Zero}{Divisibility is missing its right operand.}{$0\mid0$.}{$0\mid$.}
\KeyItem{N}{F04: Divisibility laws}{The chained display is missing its final integer.}{$a\mid b$ and $b\mid c$.}{$a\mid b\mid$.}
\KeyItem{F}{F05: Linear combinations}{Adding one need not preserve divisibility.}{}{If $a\mid b$, then $a\mid(b+1)$.}
\StageHeader{2. Read gcd, quotients, and remainders}{Identify the divisor, quotient, remainder, and required bounds before judging.}
\KeyItem{T}{F06: Greatest common divisor}{Six is their greatest positive common divisor.}{}{$\gcd(12,18)=6$.}
\KeyItem{N}{F07: Coprimality}{The first gcd argument is missing.}{$\gcd(8,15)=1$.}{$\gcd(,15)=1$.}
\KeyItem{F}{F08: Division algorithm}{The remainder condition is false.}{}{$23=5(3)+8$ with $0\le8<5$.}
\KeyItem{N}{F09: Remainders}{The modulus is missing.}{$23\bmod 5$.}{$23\bmod$.}
\KeyItem{T}{F10: Congruence}{Their difference is fifteen, divisible by five.}{}{$17\equiv_5 2$.}
\columnbreak
\StageHeader{3. Interpret modular congruence}{Translate congruence into divisibility of a difference before using a law.}
\KeyItem{F}{F11: Congruence}{The divisibility relation in the definition is reversed and altered.}{}{$a\equiv_m b\iff a\mid(m-b)$.}
\KeyItem{F}{F12: Congruence laws}{Adding $c$ to only one side need not preserve congruence.}{}{$a\equiv_m b\Rightarrow a+c\equiv_m b$.}
\KeyItem{T}{F13: Congruence laws}{Multiplying the divisible difference by $c$ preserves divisibility.}{}{$a\equiv_m b\Rightarrow ac\equiv_m bc$.}
\KeyItem{F}{F14: Cancellation}{For example, $2\cdot1\equiv_4 2\cdot3$ but $1\not\equiv_4 3$.}{}{If $ac\equiv_m bc$, then $a\equiv_m b$ without any condition on $c$.}
\KeyItem{N}{F15: Parity}{The variable $k$ is free and unassigned.}{$\exists k\in\mathbb Z,\ 14=2k$.}{$14=2k+1$.}
\StageHeader{4. Read primes, bases, and place value}{Check domains, legal digits, and place values before performing arithmetic.}
\KeyItem{T}{F16: Primes}{Its only positive divisors are one and thirteen.}{}{$13$ is prime.}
\KeyItem{N}{F17: Base notation}{Digit two is not permitted in base two.}{$(110)_2=6$.}{$(102)_2=6$.}
\KeyItem{F}{F18: Base notation}{The place-value sum is twenty-one.}{}{$(210)_3=19$.}
\KeyItem{N}{F19: Repeated division}{The quotient is missing.}{$202=2(101)+0$.}{$202=2()+0$.}
\KeyItem{F}{F20: Place value}{The right side equals $b^3$, not $b^2$.}{}{$(10)_b(10)_b=(1000)_b$ for every base $b\ge2$.}
}
\CoreForm{Divisibility}{B}{Let $a,b,c$ be arbitrary integers and let $m\ge2$ be a fixed integer. All other variables are integers when quantified or assigned. In $a\mid b$, the left integer divides the right integer. In $a\equiv_m b$, the modulus satisfies $m\ge2$. A base-$b$ numeral may use only digits $0,\ldots,b-1$.}{$a\mid b$: $a$ divides $b$; $a\nmid b$: $a$ does not divide $b$; $a\equiv b\pmod m$: $a$ and $b$ have the same remainder modulo $m$; $\gcd(a,b)$: greatest common divisor; $(d_k\cdots d_0)_b$: a base-$b$ numeral.}{
\StageHeader{1. Read divisibility from its definition}{Read $a\mid b$ as an existential equation: $b=ak$ for some integer $k$.}
\KeyItem{F}{F01: Divisibility}{Twelve is not an integer multiple of five.}{}{$5\mid12$.}
\KeyItem{F}{F02: Definition}{The left side is false while the right side has witness $k=4$.}{}{$6\mid20\iff\exists k\in\mathbb Z,\ 20=5k$.}
\KeyItem{T}{F03: Zero}{Zero equals $7(0)$.}{}{$7\mid0$.}
\KeyItem{T}{F04: Divisibility laws}{Substituting the two integer-multiple equations proves transitivity.}{}{If $a\mid b$ and $b\mid c$, then $a\mid c$.}
\KeyItem{N}{F05: Linear combinations}{The linear combination is incomplete.}{$a\mid(3b-2c)$.}{$a\mid(3b-)$.}
\StageHeader{2. Read gcd, quotients, and remainders}{Identify the divisor, quotient, remainder, and required bounds before judging.}
\KeyItem{F}{F06: Greatest common divisor}{Twelve does not divide eighteen.}{}{$\gcd(12,18)=12$.}
\KeyItem{T}{F07: Coprimality}{The two integers are coprime.}{}{$\gcd(8,15)=1$.}
\KeyItem{N}{F08: Division algorithm}{The variables $q$ and $r$ are free.}{$\exists q,r\in\mathbb Z:\ 23=5q+r$ and $0\le r<5$.}{$23=5q+r$.}
\KeyItem{T}{F09: Remainders}{Twenty-three leaves remainder three upon division by five.}{}{$23\bmod 5=3$.}
\KeyItem{F}{F10: Congruence}{Their difference is fourteen, not divisible by five.}{}{$17\equiv_5 3$.}
\columnbreak
\StageHeader{3. Interpret modular congruence}{Translate congruence into divisibility of a difference before using a law.}
\KeyItem{N}{F11: Congruence}{The modulus required by the stated notation is missing.}{$a\equiv_m b$.}{$a\equiv b$.}
\KeyItem{N}{F12: Congruence laws}{Congruence is missing its right operand.}{$a+c\equiv_m b+c$.}{$a+c\equiv_m$.}
\KeyItem{F}{F13: Congruence laws}{Multiplying only one side need not preserve congruence.}{}{$a\equiv_m b\Rightarrow ac\equiv_m b$.}
\KeyItem{N}{F14: Cancellation}{The gcd equality is missing its value.}{$\gcd(c,m)=1\Rightarrow a\equiv_m b$.}{$\gcd(c,m)=\Rightarrow a\equiv_m b$.}
\KeyItem{T}{F15: Parity}{$14=2(7)$.}{}{$14$ is even.}
\StageHeader{4. Read primes, bases, and place value}{Check domains, legal digits, and place values before performing arithmetic.}
\KeyItem{F}{F16: Primes}{Three and five are nontrivial divisors.}{}{$15$ is prime.}
\KeyItem{T}{F17: Base notation}{$1\cdot2^2+0\cdot2+1=5$.}{}{$(101)_2=5$.}
\KeyItem{N}{F18: Base notation}{Digit three is not permitted in base three.}{$(101)_3=10$.}{$(31)_3=10$.}
\KeyItem{T}{F19: Repeated division}{The equality holds and the remainder is zero.}{}{$202=2(101)+0$ is a valid first binary-conversion step.}
\KeyItem{N}{F20: Place value}{Base one is not positional notation under the stated convention.}{$(10)_2(10)_2=(100)_2$.}{$(10)_1(10)_1=(100)_1$.}
}
\CoreForm{Divisibility}{C}{Let $a,b,c$ be arbitrary integers and let $m\ge2$ be a fixed integer. All other variables are integers when quantified or assigned. In $a\mid b$, the left integer divides the right integer. In $a\equiv_m b$, the modulus satisfies $m\ge2$. A base-$b$ numeral may use only digits $0,\ldots,b-1$.}{$a\mid b$: $a$ divides $b$; $a\nmid b$: $a$ does not divide $b$; $a\equiv b\pmod m$: $a$ and $b$ have the same remainder modulo $m$; $\gcd(a,b)$: greatest common divisor; $(d_k\cdots d_0)_b$: a base-$b$ numeral.}{
\StageHeader{1. Read divisibility from its definition}{Read $a\mid b$ as an existential equation: $b=ak$ for some integer $k$.}
\KeyItem{N}{F01: Divisibility}{Divisibility is missing its right operand.}{$3\mid12$.}{$3\mid$.}
\KeyItem{N}{F02: Definition}{The variable $k$ is free and has not been assigned.}{$\exists k\in\mathbb Z,\ 20=4k$.}{$20=4k$.}
\KeyItem{F}{F03: Zero}{No integer multiplied by zero equals seven.}{}{$0\mid7$.}
\KeyItem{F}{F04: Divisibility laws}{Divisibility is not generally symmetric.}{}{If $a\mid b$, then $b\mid a$.}
\KeyItem{T}{F05: Linear combinations}{Any integer linear combination remains divisible by $a$.}{}{If $a\mid b$ and $a\mid c$, then $a\mid(3b-2c)$.}
\StageHeader{2. Read gcd, quotients, and remainders}{Identify the divisor, quotient, remainder, and required bounds before judging.}
\KeyItem{N}{F06: Greatest common divisor}{The binary gcd operation is missing an argument.}{$\gcd(12,18)$.}{$\gcd(12,)$.}
\KeyItem{F}{F07: Coprimality}{Their gcd is four.}{}{$\gcd(8,12)=1$.}
\KeyItem{T}{F08: Division algorithm}{This is the valid quotient-remainder representation.}{}{$23=5(4)+3$ with $0\le3<5$.}
\KeyItem{F}{F09: Remainders}{A remainder modulo five must be below five.}{}{$23\bmod 5=5$.}
\KeyItem{N}{F10: Congruence}{Congruence is missing its right operand.}{$17\equiv_5 2$.}{$17\equiv_5$.}
\columnbreak
\StageHeader{3. Interpret modular congruence}{Translate congruence into divisibility of a difference before using a law.}
\KeyItem{T}{F11: Congruence}{This is the definition of congruence modulo $m$.}{}{$a\equiv_m b\iff m\mid(a-b)$.}
\KeyItem{T}{F12: Congruence laws}{Adding the same integer preserves the difference.}{}{$a\equiv_m b\Rightarrow a+c\equiv_m b+c$.}
\KeyItem{N}{F13: Congruence laws}{The conjunction is incomplete.}{$ac\equiv_m bc$.}{$ac\equiv_m bc\land$.}
\KeyItem{T}{F14: Cancellation}{A factor coprime to the modulus can be cancelled.}{}{If $ac\equiv_m bc$ and $\gcd(c,m)=1$, then $a\equiv_m b$.}
\KeyItem{F}{F15: Parity}{Fourteen has the even form, not the odd form.}{}{$14$ is odd.}
\StageHeader{4. Read primes, bases, and place value}{Check domains, legal digits, and place values before performing arithmetic.}
\KeyItem{N}{F16: Primes}{The explanation is syntactically incomplete.}{$13$ is prime because its only positive divisors are $1$ and $13$.}{$1$ is prime because.}
\KeyItem{F}{F17: Base notation}{The place-value sum is five.}{}{$(101)_2=6$.}
\KeyItem{T}{F18: Base notation}{$2\cdot3^2+1\cdot3=21$.}{}{$(210)_3=21$.}
\KeyItem{F}{F19: Repeated division}{The remainder must be zero or one.}{}{$202=2(100)+2$ is a valid first binary-conversion step.}
\KeyItem{T}{F20: Place value}{Both sides equal $b^2$.}{}{$(10)_b(10)_b=(100)_b$ for every base $b\ge2$.}
}
\CoreForm{Induction}{A}{$P(n)$ is a proposition indexed by integers $n\ge1$. In an induction step, $k$ is a fixed but arbitrary integer in the stated range. Also let $a_1=2$ and $a_{n+1}=a_n+3$ for $n\ge1$.}{$P(k)$: the statement at index $k$; $P(k)\to P(k+1)$: the induction step; $a_n$: the $n$th sequence term; $\sum$: add the indexed terms; strong induction may use all earlier established cases.}{
\StageHeader{1. Read the parts of an induction proof}{Label the base case, hypothesis, and step; keep assumptions separate from goals.}
\KeyItem{T}{F01: Base case}{The base must match the theorem's first index.}{}{For a theorem beginning at $n=1$, the base case is $P(1)$.}
\KeyItem{F}{F02: Induction hypothesis}{That is circular.}{}{The induction hypothesis assumes $P(k+1)$ in order to prove $P(k+1)$.}
\KeyItem{F}{F03: Induction step}{That does not advance to a new case.}{}{The induction step proves only $P(k)\to P(k)$.}
\KeyItem{T}{F04: Induction principle}{This is the principle of ordinary induction.}{}{$P(1)$ and $\forall k\ge1\,[P(k)\to P(k+1)]$ imply $\forall n\ge1\,P(n)$.}
\KeyItem{T}{F05: Induction logic}{A finite list lacks the propagation argument.}{}{Checking several initial cases does not prove every later case.}
\StageHeader{2. Track hypotheses and base coverage}{Track exactly which earlier cases are available and which starting cases are required.}
\KeyItem{N}{F06: Arbitrary index}{The equality is incomplete.}{Let $k\ge1$ be arbitrary.}{Let $k=$ be arbitrary.}
\KeyItem{F}{F07: Strong induction}{That is circular.}{}{Strong induction assumes $P(k+1)$ to prove $P(k+1)$.}
\KeyItem{N}{F08: Strong induction}{The conditional has no consequent.}{$P(1)\land\cdots\land P(k)\to P(k+1)$.}{$P(1)\land\cdots\land P(k)\to$.}
\KeyItem{T}{F09: Base coverage}{The induction step may need both earlier values.}{}{A recurrence using two preceding terms may require two initial cases.}
\KeyItem{N}{F10: Recurrences}{The right side is incomplete.}{$a_{n+1}=a_n+3$.}{$a_{n+1}=a_n+$.}
\columnbreak
\StageHeader{3. Read recurrences and closed forms}{Substitute an index explicitly before comparing a recurrence with a closed form.}
\KeyItem{N}{F11: Initial conditions}{The initial-value equality is incomplete.}{$a_1=2$.}{$a_1=$.}
\KeyItem{T}{F12: Sequence values}{The sequence is $2,5,8,11,\ldots$.}{}{$a_4=11$.}
\KeyItem{T}{F13: Closed forms}{It matches $a_1=2$ and increases by three.}{}{$a_n=3n-1$ for every $n\ge1$.}
\KeyItem{F}{F14: Induction algebra}{Substitution requires an established equality for that expression.}{}{The induction hypothesis may replace a different expression that it does not mention.}
\KeyItem{N}{F15: Summation identity}{The variable $n$ is free in this standalone formula.}{$\forall n\ge1,\ 1+\cdots+n=\frac{n(n+1)}2$.}{$1+\cdots+n=\frac{n(n+1)}2$.}
\StageHeader{4. Audit the algebra and proof logic}{Check the claimed range and verify that each algebraic line follows from the hypothesis.}
\KeyItem{F}{F16: Inequalities}{The claim already fails at $n=1$.}{}{$2^n<n$ for every integer $n\ge1$.}
\KeyItem{N}{F17: Implication use}{The conditional premise is incomplete.}{$P(k),\ P(k)\to P(k+1)\ \therefore P(k+1)$.}{$P(k),\ P(k)\to\ \therefore P(k+1)$.}
\KeyItem{F}{F18: Circularity}{That assumes the desired next case.}{}{Assuming $P(k+1)$ is allowed while proving $P(k+1)$.}
\KeyItem{F}{F19: Starting index}{That base lies outside the claimed range.}{}{A theorem stated for $n\ge4$ is proved by checking only $n=1$.}
\KeyItem{N}{F20: Proposition/term}{Arithmetic addition is not defined for a proposition.}{$a_k+1$ or $P(k+1)$.}{$P(k)+1$.}
}
\CoreForm{Induction}{B}{$P(n)$ is a proposition indexed by integers $n\ge1$. In an induction step, $k$ is a fixed but arbitrary integer in the stated range. Also let $a_1=2$ and $a_{n+1}=a_n+3$ for $n\ge1$.}{$P(k)$: the statement at index $k$; $P(k)\to P(k+1)$: the induction step; $a_n$: the $n$th sequence term; $\sum$: add the indexed terms; strong induction may use all earlier established cases.}{
\StageHeader{1. Read the parts of an induction proof}{Label the base case, hypothesis, and step; keep assumptions separate from goals.}
\KeyItem{F}{F01: Base case}{Zero lies outside the stated starting range.}{}{For a theorem beginning at $n=1$, the base case is necessarily $P(0)$.}
\KeyItem{N}{F02: Induction hypothesis}{The conjunction is incomplete.}{Assume $P(k)$.}{Assume $P(k)\land$.}
\KeyItem{N}{F03: Induction step}{The predicate in the consequent has no argument.}{$P(k)\to P(k+1)$.}{$P(k)\to P()$.}
\KeyItem{F}{F04: Induction principle}{A base case does not propagate without an induction step.}{}{$P(1)$ alone implies $\forall n\ge1\,P(n)$.}
\KeyItem{F}{F05: Induction logic}{Later cases remain unproved.}{}{Checking $P(1),P(2),P(3)$ proves $\forall n\ge1\,P(n)$.}
\StageHeader{2. Track hypotheses and base coverage}{Track exactly which earlier cases are available and which starting cases are required.}
\KeyItem{T}{F06: Arbitrary index}{A special choice would establish only a special transition.}{}{The induction step must work for an arbitrary allowed value of $k$.}
\KeyItem{N}{F07: Strong induction}{The list of assumed cases has no endpoint.}{Assume $P(1),\ldots,P(k)$.}{Assume $P(1),\ldots,$.}
\KeyItem{T}{F08: Strong induction}{They are equivalent proof principles.}{}{Ordinary and strong induction can prove the same class of statements.}
\KeyItem{F}{F09: Base coverage}{The number of required starting cases depends on the statement and step.}{}{Every induction proof has exactly one base case.}
\KeyItem{T}{F10: Recurrences}{That is the defining role of a recurrence.}{}{A recurrence defines later terms using earlier terms.}
\columnbreak
\StageHeader{3. Read recurrences and closed forms}{Substitute an index explicitly before comparing a recurrence with a closed form.}
\KeyItem{T}{F11: Initial conditions}{The recurrence and starting value determine each later term.}{}{The recurrence $a_{n+1}=a_n+3$ and $a_1=2$ determine a unique sequence.}
\KeyItem{F}{F12: Sequence values}{Eight is the third term.}{}{$a_4=8$.}
\KeyItem{F}{F13: Closed forms}{At $n=1$ it gives four instead of two.}{}{$a_n=3n+1$ for every $n\ge1$.}
\KeyItem{N}{F14: Induction algebra}{The equality is missing its right-hand expression.}{$1+\cdots+k=\frac{k(k+1)}2$.}{$1+\cdots+k=$.}
\KeyItem{T}{F15: Summation identity}{This is the standard arithmetic-series identity.}{}{$\forall n\ge1,\ 1+\cdots+n=\frac{n(n+1)}2$.}
\StageHeader{4. Audit the algebra and proof logic}{Check the claimed range and verify that each algebraic line follows from the hypothesis.}
\KeyItem{N}{F16: Inequalities}{The inequality is missing its right-hand expression.}{$2^n\ge n+1$.}{$2^n\ge$.}
\KeyItem{T}{F17: Implication use}{This is modus ponens inside the induction proof.}{}{From $P(k)$ and $P(k)\to P(k+1)$, one may infer $P(k+1)$.}
\KeyItem{N}{F18: Circularity}{The justification after because is missing.}{Assume $P(k)$ as the induction hypothesis.}{Assume $P(k+1)$ because.}
\KeyItem{N}{F19: Starting index}{The lower bound is missing.}{$n\ge4$.}{$n\ge$.}
\KeyItem{T}{F20: Proposition/term}{Excluded middle makes this proposition true.}{}{$P(1)\lor\neg P(1)$.}
}
\CoreForm{Induction}{C}{$P(n)$ is a proposition indexed by integers $n\ge1$. In an induction step, $k$ is a fixed but arbitrary integer in the stated range. Also let $a_1=2$ and $a_{n+1}=a_n+3$ for $n\ge1$.}{$P(k)$: the statement at index $k$; $P(k)\to P(k+1)$: the induction step; $a_n$: the $n$th sequence term; $\sum$: add the indexed terms; strong induction may use all earlier established cases.}{
\StageHeader{1. Read the parts of an induction proof}{Label the base case, hypothesis, and step; keep assumptions separate from goals.}
\KeyItem{N}{F01: Base case}{The unary predicate has no argument.}{The base case is $P(1)$.}{The base case is $P()$.}
\KeyItem{T}{F02: Induction hypothesis}{It supplies the current case for proving the next one.}{}{The induction hypothesis assumes $P(k)$ for an arbitrary $k\ge1$.}
\KeyItem{T}{F03: Induction step}{It establishes that truth passes to the next index.}{}{The induction step proves $P(k)\to P(k+1)$.}
\KeyItem{N}{F04: Induction principle}{The universal quantifier over $k$ has no formula in its scope.}{$P(1)\land\forall k\ge1\,[P(k)\to P(k+1)]$.}{$P(1)\land\forall k\ge1\Rightarrow\forall n\,P(n)$.}
\KeyItem{N}{F05: Induction logic}{The displayed argument has no conclusion.}{$P(1),P(2),P(3)\therefore P(1)\land P(2)\land P(3)$.}{$P(1),P(2),P(3)\therefore$.}
\StageHeader{2. Track hypotheses and base coverage}{Track exactly which earlier cases are available and which starting cases are required.}
\KeyItem{F}{F06: Arbitrary index}{It verifies only one transition.}{}{Proving $P(7)\to P(8)$ establishes the complete induction step for every $k$.}
\KeyItem{T}{F07: Strong induction}{It makes all earlier cases available.}{}{Strong induction may assume $P(1),\ldots,P(k)$ to prove $P(k+1)$.}
\KeyItem{F}{F08: Strong induction}{It is a convenient formulation, not a stronger logic.}{}{Strong induction is valid only when ordinary induction fails.}
\KeyItem{N}{F09: Base coverage}{The list ends after a comma.}{Base cases: $P(1),P(2)$.}{Base cases: $P(1),$.}
\KeyItem{F}{F10: Recurrences}{Different initial values can produce different sequences.}{}{A recurrence alone always determines a unique sequence without initial values.}
\columnbreak
\StageHeader{3. Read recurrences and closed forms}{Substitute an index explicitly before comparing a recurrence with a closed form.}
\KeyItem{F}{F11: Initial conditions}{Each possible starting value gives a different sequence.}{}{The recurrence $a_{n+1}=a_n+3$ determines one unique sequence without $a_1$.}
\KeyItem{N}{F12: Sequence values}{This is an incomplete arithmetic term, not a proposition.}{$a_4=11$.}{$a_{4}+$.}
\KeyItem{N}{F13: Closed forms}{The variable $n$ is free in this standalone formula.}{$\forall n\ge1,\ a_n=3n-1$.}{$a_n=3n-1$.}
\KeyItem{T}{F14: Induction algebra}{That is the hypothesis's permitted use.}{}{In a sum proof, the induction hypothesis may replace the partial sum through $k$.}
\KeyItem{F}{F15: Summation identity}{At $n=1$ the right side is zero.}{}{$\forall n\ge1,\ 1+\cdots+n=\frac{n(n-1)}2$.}
\StageHeader{4. Audit the algebra and proof logic}{Check the claimed range and verify that each algebraic line follows from the hypothesis.}
\KeyItem{T}{F16: Inequalities}{The statement has a valid base and induction step.}{}{$2^n\ge n+1$ for every integer $n\ge0$.}
\KeyItem{F}{F17: Implication use}{That is affirming the consequent.}{}{From $P(k+1)$ and $P(k)\to P(k+1)$, one may infer $P(k)$.}
\KeyItem{T}{F18: Circularity}{It is the antecedent of the conditional being proved.}{}{Assuming $P(k)$ is allowed while proving $P(k)\to P(k+1)$.}
\KeyItem{T}{F19: Starting index}{The base case must start where the theorem starts.}{}{A theorem stated for $n\ge4$ should begin its induction at $n=4$.}
\KeyItem{F}{F20: Proposition/term}{A proposition and its negation cannot both hold.}{}{$P(1)\land\neg P(1)$.}
}
\CoreForm{Counting}{A}{Let $X=\{a,b\}$, $Y=\{b,c,d\}$, $n=5$, and $r=2$. For probability, a fair die has sample space $\Omega=\{1,2,3,4,5,6\}$ with events $E=\{2,4,6\}$ and $F=\{4,5,6\}$. Let $Z$ be the number shown.}{$n!$: factorial; $P(n,r)$: ordered selections; $\binom nr$: unordered selections; $|X|$: size of $X$; $\Pr(E)$: probability of event $E$; $E^c$: not $E$; $\Pr(E\mid F)$: probability of $E$ given $F$; $\mathbb E[Z]$: expected value.}{
\StageHeader{1. Choose and read a counting rule}{Decide whether stages are combined by addition or multiplication before calculating.}
\KeyItem{N}{F01: Sum rule}{The sum and equality are incomplete.}{$3+4=7$.}{$3+=$.}
\KeyItem{F}{F02: Product rule}{Seven would add alternatives rather than combine stages.}{}{A two-stage process with $3$ choices and then $4$ choices has $7$ outcomes.}
\KeyItem{F}{F03: Inclusion-exclusion}{The shared element $b$ must not be counted twice.}{}{$|X\cup Y|=5$.}
\KeyItem{N}{F04: Pigeonhole principle}{The argument has no conclusion.}{$7$ objects in $6$ boxes $\therefore$ a repeated box.}{$7\to6\therefore$.}
\KeyItem{N}{F05: Factorials}{Factorial is not defined for the set $X$ under these conventions.}{$|X|!=2$.}{$X!$.}
\StageHeader{2. Interpret counting notation and models}{Ask whether order and repetition matter, then match the notation to that model.}
\KeyItem{T}{F06: Permutations}{There are $5\cdot4=20$ ordered selections.}{}{$P(5,2)=20$.}
\KeyItem{F}{F07: Combinations}{Twenty counts ordered selections.}{}{$\binom52=20$.}
\KeyItem{T}{F08: Combinations}{Choosing two to include is equivalent to choosing three to exclude.}{}{$\binom52=\binom53$.}
\KeyItem{T}{F09: Strings}{Each of five positions has two choices.}{}{There are $2^5=32$ bit strings of length $5$.}
\KeyItem{F}{F10: Repetition}{The permutation formula forbids repetition.}{}{Five choices made twice with repetition produce $P(5,2)$ sequences.}
\columnbreak
\StageHeader{3. Read events and probability rules}{Identify the sample space and event type before applying a probability rule.}
\KeyItem{N}{F11: Repeated objects}{The denominator is incomplete.}{$\frac{5!}{2!\,2!}=30$.}{$\frac{5!}{2!\,}$.}
\KeyItem{T}{F12: Order}{Committee membership is unordered.}{}{Selecting a two-person committee from five people is counted by $\binom52$.}
\KeyItem{F}{F13: Probability type}{Seven is not a possible die outcome.}{}{$7\in\Omega$.}
\KeyItem{F}{F14: Classical probability}{The favorable-to-total ratio is $3/6$.}{}{$\Pr(E)=\frac13$.}
\KeyItem{T}{F15: Complements}{An event and its complement partition the sample space.}{}{$\Pr(E^c)=1-\Pr(E)$.}
\StageHeader{4. Read dependence, conditioning, and expectation}{Check conditioning, dependence, and the type of the object inside each operator.}
\KeyItem{N}{F16: Addition rule}{Union is missing its second event.}{$\Pr(E\cup F)$.}{$\Pr(E\cup)=1$.}
\KeyItem{N}{F17: Disjointness}{Intersection is missing its second event.}{$E\cap F$ is nonempty.}{$E\cap$ is empty.}
\KeyItem{T}{F18: Independence}{The intersection is $\{4,6\}$, two of six outcomes.}{}{$\Pr(E\cap F)=\frac13$.}
\KeyItem{N}{F19: Conditional probability}{Conditional probability is undefined when the conditioning event has probability zero.}{$\Pr(E\mid F)$.}{$\Pr(E\mid\varnothing)$.}
\KeyItem{F}{F20: Expected value}{The arithmetic mean is three and one half.}{}{$\mathbb E[Z]=3$.}
}
\CoreForm{Counting}{B}{Let $X=\{a,b\}$, $Y=\{b,c,d\}$, $n=5$, and $r=2$. For probability, a fair die has sample space $\Omega=\{1,2,3,4,5,6\}$ with events $E=\{2,4,6\}$ and $F=\{4,5,6\}$. Let $Z$ be the number shown.}{$n!$: factorial; $P(n,r)$: ordered selections; $\binom nr$: unordered selections; $|X|$: size of $X$; $\Pr(E)$: probability of event $E$; $E^c$: not $E$; $\Pr(E\mid F)$: probability of $E$ given $F$; $\mathbb E[Z]$: expected value.}{
\StageHeader{1. Choose and read a counting rule}{Decide whether stages are combined by addition or multiplication before calculating.}
\KeyItem{T}{F01: Sum rule}{Disjoint alternatives are added.}{}{If two disjoint cases have $3$ and $4$ outcomes, together they have $7$ outcomes.}
\KeyItem{N}{F02: Product rule}{Multiplication is missing its second factor.}{$3\cdot4=12$.}{$3\cdot$.}
\KeyItem{N}{F03: Inclusion-exclusion}{Union is missing its second set.}{$|X\cup Y|=4$.}{$|X\cup|=4$.}
\KeyItem{T}{F04: Pigeonhole principle}{There are more objects than boxes.}{}{Placing $7$ objects into $6$ boxes forces some box to receive at least two objects.}
\KeyItem{T}{F05: Factorials}{$5\cdot4\cdot3\cdot2\cdot1=120$.}{}{$5!=120$.}
\StageHeader{2. Interpret counting notation and models}{Ask whether order and repetition matter, then match the notation to that model.}
\KeyItem{F}{F06: Permutations}{Ten counts unordered selections.}{}{$P(5,2)=10$.}
\KeyItem{N}{F07: Combinations}{The population size is missing.}{$\binom52=10$.}{$\binom{\,}{2}=10$.}
\KeyItem{F}{F08: Combinations}{The values are ten and five.}{}{$\binom52=\binom54$.}
\KeyItem{F}{F09: Strings}{Bits may repeat, so permutations are not the correct model.}{}{There are $5!=120$ bit strings of length $5$.}
\KeyItem{N}{F10: Repetition}{The box is an unfilled sequence-length placeholder.}{$5^2$ sequences.}{$5^{\Box}$ sequences.}
\columnbreak
\StageHeader{3. Read events and probability rules}{Identify the sample space and event type before applying a probability rule.}
\KeyItem{T}{F11: Repeated objects}{The two Ls and two Es are indistinguishable within their groups.}{}{The word LEVEL has $\frac{5!}{2!\,2!}=30$ distinct arrangements.}
\KeyItem{F}{F12: Order}{The offices make the selection ordered.}{}{Assigning president and secretary from five people is counted by $\binom52$.}
\KeyItem{N}{F13: Probability type}{Membership requires a set on the right.}{$4\in E$.}{$E\in 4$.}
\KeyItem{N}{F14: Classical probability}{Probability is applied to events; the outcome should be written as a singleton event.}{$\Pr(\{2\})=\frac16$.}{$\Pr(2)=\frac16$.}
\KeyItem{F}{F15: Complements}{The complement rule subtracts from one.}{}{$\Pr(E^c)=\Pr(E)+1$.}
\StageHeader{4. Read dependence, conditioning, and expectation}{Check conditioning, dependence, and the type of the object inside each operator.}
\KeyItem{T}{F16: Addition rule}{The intersection corrects double counting.}{}{$\Pr(E\cup F)=\Pr(E)+\Pr(F)-\Pr(E\cap F)$.}
\KeyItem{T}{F17: Disjointness}{Their intersection is empty.}{}{$\{1\}$ and $\{2\}$ are disjoint events.}
\KeyItem{F}{F18: Independence}{$\Pr(E\cap F)=1/3$ but $\Pr(E)\Pr(F)=1/4$.}{}{$E$ and $F$ are independent.}
\KeyItem{T}{F19: Conditional probability}{Two of the three outcomes in $F$ are even.}{}{$\Pr(E\mid F)=\frac23$.}
\KeyItem{N}{F20: Expected value}{The expectation operator is defined here for the random variable $Z$, not the event $E$.}{$\mathbb E[Z]$.}{$\mathbb E[E]$.}
}
\CoreForm{Counting}{C}{Let $X=\{a,b\}$, $Y=\{b,c,d\}$, $n=5$, and $r=2$. For probability, a fair die has sample space $\Omega=\{1,2,3,4,5,6\}$ with events $E=\{2,4,6\}$ and $F=\{4,5,6\}$. Let $Z$ be the number shown.}{$n!$: factorial; $P(n,r)$: ordered selections; $\binom nr$: unordered selections; $|X|$: size of $X$; $\Pr(E)$: probability of event $E$; $E^c$: not $E$; $\Pr(E\mid F)$: probability of $E$ given $F$; $\mathbb E[Z]$: expected value.}{
\StageHeader{1. Choose and read a counting rule}{Decide whether stages are combined by addition or multiplication before calculating.}
\KeyItem{F}{F01: Sum rule}{Multiplication is for successive choices, not alternative cases.}{}{If two disjoint cases have $3$ and $4$ outcomes, together they have $12$ outcomes.}
\KeyItem{T}{F02: Product rule}{Successive independent choice counts multiply.}{}{A two-stage process with $3$ choices and then $4$ choices has $12$ outcomes.}
\KeyItem{T}{F03: Inclusion-exclusion}{The union is $\{a,b,c,d\}$.}{}{$|X\cup Y|=4$.}
\KeyItem{F}{F04: Pigeonhole principle}{One object can be placed in each box.}{}{Placing $6$ objects into $6$ boxes forces some box to receive at least two objects.}
\KeyItem{F}{F05: Factorials}{Factorial is not repeated addition or squaring.}{}{$5!=25$.}
\StageHeader{2. Interpret counting notation and models}{Ask whether order and repetition matter, then match the notation to that model.}
\KeyItem{N}{F06: Permutations}{The number selected is missing.}{$P(5,2)=20$.}{$P(5,)=20$.}
\KeyItem{T}{F07: Combinations}{There are ten unordered pairs from five objects.}{}{$\binom52=10$.}
\KeyItem{N}{F08: Combinations}{The lower argument must be an integer, not the set $X$.}{$\binom52=10$.}{$\binom5X=10$.}
\KeyItem{N}{F09: Strings}{The box is an unfilled exponent placeholder, not a number.}{$2^5=32$.}{$2^{\Box}=32$.}
\KeyItem{T}{F10: Repetition}{Each of the two positions has five choices.}{}{Five choices made twice with repetition produce $5^2=25$ sequences.}
\columnbreak
\StageHeader{3. Read events and probability rules}{Identify the sample space and event type before applying a probability rule.}
\KeyItem{F}{F11: Repeated objects}{That count treats repeated letters as distinct.}{}{The word LEVEL has $5!=120$ distinct arrangements.}
\KeyItem{N}{F12: Order}{The requested committee size is missing.}{$\binom52$ committees.}{committee$(5,)$.}
\KeyItem{T}{F13: Probability type}{An event is a subset of the sample space.}{}{$E\subseteq\Omega$.}
\KeyItem{T}{F14: Classical probability}{Three of six equally likely outcomes are even.}{}{$\Pr(E)=\frac12$.}
\KeyItem{N}{F15: Complements}{The event being complemented is missing.}{$\Pr(E^c)=1-\Pr(E)$.}{$\Pr({}^c)=1$.}
\StageHeader{4. Read dependence, conditioning, and expectation}{Check conditioning, dependence, and the type of the object inside each operator.}
\KeyItem{F}{F16: Addition rule}{The events overlap at outcomes four and six.}{}{$\Pr(E\cup F)=\Pr(E)+\Pr(F)$.}
\KeyItem{F}{F17: Disjointness}{They share outcomes four and six.}{}{$E$ and $F$ are disjoint events.}
\KeyItem{N}{F18: Independence}{The second probability has no event argument.}{$\Pr(E)\Pr(F)=\Pr(E\cap F)$.}{$\Pr(E)\Pr()=\Pr(E\cap F)$.}
\KeyItem{F}{F19: Conditional probability}{Conditioning restricts the denominator to $F$.}{}{$\Pr(E\mid F)=\frac12$.}
\KeyItem{T}{F20: Expected value}{The mean of the six equally likely die values is $21/6$.}{}{$\mathbb E[Z]=3.5$.}
}

\end{document}
